% ----------Introduction----------
\section{Motivation}

% General about language learning (LL) -> importance -> why has grown -> mobile-assisted
The scale of foreign language acquisition is an unprecedented phenomenon in contemporary applied linguistics. In our globalised society, over 1.5 billion individuals are actively studying a foreign language \cite{kentstate2025language}, seeking to facilitate cross-cultural communication, academic mobility, and access to global scientific literature. This widespread interest in language learning in conjunction with the advent of technological advancement (e.g., smartphones, AI) has shifted the paradigm of second language acquisition (SLA) away from traditional classroom-bound instruction toward self-directed, digital learning environments. Currently, independent, mobile-assisted platforms represent the most widespread medium for independent language study, enabling millions of learners to engage with language materials on demand \cite{mordor2026language}.

% Mobile-assisted LL problems 
Despite the ubiquity of mobile-first access and digital resources, contemporary Mobile-Assisted Language Learning (MALL) faces critical pedagogical efficacy bottlenecks. While self-learning mobile applications have become the standard interface for independent L2 (second language) acquisition \cite{mordor2026language}, their long-term educational utility is severely compromised by learner attrition. Industrial Analysis of self-directed language learning applications reveals a severe pedagogical retention deficit, with first-week user attrition rates exceeding 60\% across many major platforms \cite{mordor2026language}.

% No time for long-form classroom study -> micro-learning
Furthermore, modern lifestyle demands and cognitive workloads compress the time available for structured, long-form study \cite{talentlms2026}. Traditional, exhaustive learning structures fail to integrate into the daily routines of adult learners, as evidenced by empirical findings indicating that 53\% of active adult learners and half of educational directors identify excessive daily workloads as the primary barrier to ongoing instruction \cite{talentlms2026}. To circumvent these temporal constraints, learner behaviors have increasingly shifted toward autonomous, bit-sized micro-learning formats.

% micro-learning ->  Lack , limitations and deficits with the current MALL platforms -> leverage attention habits of modern learners
% The most important paragraph
% Added question for hook 
While modern mobile applications have successfully popularised the micro-learning format, current digital implementations remain heavily constrained by structural and pedagogical bottlenecks. As analyzed in depth within the literature review (Chapter~\ref{chapter:related_work}), existing MALL platforms frequently rely on uniform, static curricula that fail to adapt to individual learner demographics, lack authentic multimedia integration, and reduce gamification to repetitive, decontextualized translation drills---a phenomenon termed ``shallow gamification.'' Consequently, these platforms fail to support advanced proficiency development or sustain long-term cognitive engagement. This systemic deficit shifts the scientific inquiry: \textit{how can mobile language learning frameworks transcend these repetitive, drill-based limitations to deliver adaptive, culturally authentic, and cognitively immersive media that aligns with the organic attention habits of modern learners?} Addressing this question requires investigating informal digital environments where contemporary learners willingly and habitually direct their attention.


% ----------BACKGROUND----------
\section{Background}

    To address these completion and engagement deficits, modern Mobile-Assisted Language Learning (MALL) must align with the media consumption realities of contemporary learners. This thesis explores the integration of two prominent digital phenomena: short-form video and hypercasual game mechanics.

    \subsection{Short-form Video}
    Short-form video is a media format native to contemporary social media platforms, characterized by its brevity---typically ranging from 15 seconds to 3 minutes---and optimized for rapid, algorithmically driven consumption.

    Short-form video has become the dominant interface of digital attention across social networks. Statistically, YouTube Shorts has achieved over 200 billion daily views \cite{mohan2026letter}, 
    while Meta Platforms announced: ``Video is a particular bright spot, with video time spent on Instagram up more than 30\% since last year and 5\% more time spent on Facebook'' \cite{zuckerberg2025q3}.
    TikTok's global active advertising audience has reached 1.59 billion adults, representing 27.5\% of the global population aged 18 and older \cite{kepios2025tiktok}. 

    This media format captures a disproportionate share of mobile attention. Globally, TikTok Android users average 1 hour and 37 minutes per day in the application, compared with 1 hour and 25 minutes for YouTube, 1 hour and 13 minutes for Instagram, and 1 hour and 7 minutes for Facebook \cite{similarweb2025}.
    
    This behavioral pattern is heavily concentrated among younger cohorts: 75\% of US teens visit YouTube daily, and 61\% engage with TikTok daily \cite{pew2025teens}. Among US adults, daily TikTok usage exhibits a sharp ten-to-one generational gap, captured by 50\% daily usage among 18-to-29-year-olds compared to just 5\% among those aged 65 and older \cite{pew2025social}. 

    Pedagogically, the cognitive appeal of short-form video lies in its brevity and high visual density. While platforms keep raising video length caps (with TikTok allowing up to 10 minutes and Reels up to 3 minutes) \cite{postplanify2026}, consumer behavior remains heavily anchored around rapid completion. Wistia's benchmark analysis of 13 million videos reveals that videos under 60 seconds hold a 52\% average engagement rate, meaning viewers watch more than half the runtime before scrolling \cite{wistia2026}. In contrast, longer videos see completion rates decay exponentially \cite{wistia2026, vidyard2024}. Short-form videos therefore represent a highly engaging, authentic delivery mechanism for Second Language Acquisition (SLA).

   \subsection{Hypercasual and Hybrid-Casual Games} \label{sub:hypercasual_and_hybrid_casual_games}

   % Hypercasual games Definition/Designation -> core chame mechanics -> hybrid-casual games
    Hypercasual video games are defined as highly accessible, mobile-native software products characterized by minimalist user interfaces, instantaneous ``tap-to-play'' mechanics, and extremely short core gameplay loops \cite{Zelevskaya2022thesis}. Unlike casual games, which require some level of long-term planning, tutorial acquisition, or system familiarity, hyper-casual games demand no pre-existing cognitive schema or learning curve from the user, relying on immediately intuitive controls (core game mechanics) such as direction and drawing, matching, puzzle, rising and falling, swerving, swiping, tapping/timing and word and board \cite{Zelevskaya2022thesis, pizzo2023hypercasual}. In recent years, this genre has evolved into ``hybrid-casual'' gaming, which retains the low-barrier entry and core mechanics of hypercasual games while strategically incorporating mid-core elements, such as progressive meta-features, item collection, and light narrative layer systems \cite{pizzo2023hypercasual}.

    % The addictiveness, exponential growth
    The exceptional appeal and potential addictiveness of these games are central to their design and rapid market expansion \cite{pizzo2023hypercasual}. Reflecting this widespread appeal, the hyper-casual and hybrid-casual genres have experienced exponential audience expansion globally, emerging as a dominant leisure activity on both major mobile platforms \cite{businessresearch2026Hypercasualgaming}. The addictiveness of these games are historically highlighted. Flappy Bird, which was voluntarily withdrawn from download by its creator due to concerns over its addictive and habit-forming impact on players \cite{pizzo2023hypercasual}.

    % wide demographics
    Furthermore, because they are playable primarily on ubiquitous mobile devices and utilize intuitive user interfaces, hypercasual games reach a broad demographic spectrum. Notably, they appeal strongly to populations that historically do not identify as gamers or have been marginalized by the ``male gamer stereotype'' \cite{Paassen2017WhatIA}. These include female players---who comprise 63\% of the mobile gaming market \cite{anderton2020hypercasual}---as well as older adults (aged 65+) seeking cognitive maintenance, mental health preservation, and familial bonding \cite{celatti2022seniorgamers}.

    % People only kill time by playing these games, they don't learn anything 
    Despite this high volume of daily engagement, players primarily consume these games as simple digital leisure pursuits to pass time, alleviate boredom, or fill transitional moments during daily routines and commutes \cite{Zelevskaya2022thesis, pizzo2023hypercasual}. The highly engaging, low-friction mechanics of hypercasual and hybrid-casual games represent a powerful, untapped attention-delivery system. While their mass appeal stems precisely from this cognitively undemanding activity, the same properties suggest a different use. Hypercasual games can already adapt their difficulty to the individual player's observed performance, which improves retention and replay value \cite{yang2020hyper} (Section~\ref{subsubsec:dda}); the same adaptivity could be directed at what the player knows rather than only at how well they play. This raises the question: can the attention that players habitually spend passing time in hypercasual games be redirected toward language learning?
    

    
% ----------PROBLEM-STATEMENT---------- 
\section{Problem Statement and Research Questions}
Short-form video and hypercasual games already capture a large and habitual share of learners' attention, yet neither is designed to teach a language. This thesis asks whether that attention can be redirected toward language learning, and how to make the result effective rather than merely engaging. Doing so raises three design problems. First, each learner's vocabulary must be tracked word by word, so that the system knows which words the learner has met and which are due for review. Second, a recommendation algorithm must use this word-level state to select reels that are both comprehensible and lexically challenging. Third, each game level must be generated from the same word-level state, so that a round of play is also a round of practice.

The central research problem is the absence of a language learning system that simultaneously leverages (i) the authentic contextual richness of user-generated short-form video and (ii) gamified reinforcement adapted to the individual learner's vocabulary state.

Research questions:

\begin{enumerate}
    \item[\textbf{RQ1.}] How can a mobile application integrate authentic short-form video content with a word-level vocabulary tracking model to support second language acquisition?
    \item[\textbf{RQ2.}] How should a content recommendation engine be designed to select reels that are both comprehensible and lexically challenging for a given learner, based on forgetting curves?
    \item[\textbf{RQ3.}] To what extent can hypercasual game mechanics, when parameterised by the learner's vocabulary state, serve as effective spaced-repetition exercises within a language learning application?
\end{enumerate}

A natural objection is that a learner could already acquire vocabulary by watching reels on existing social media platforms. Section~\ref{section:reels_in_sla} shows why this is not sufficient: passive viewing delivers comprehensible input, but provides no tracked vocabulary state, no scheduled review, and no active recall.

This thesis proposes \textbf{Glosy}, a cross-platform mobile application that addresses the identified gaps by embedding vocabulary acquisition within short-form authentic video content and reinforcing it through personalised hypercasual games. The system tracks each learner's vocabulary state at the word level, uses that state to select contextually appropriate reels, and builds the words of every game round from the same vocabulary model.


% ----------OBJECTIVES, SCOPE AND APPROACH----------
\section{Objectives, Scope, and Approach}
\label{sec:objectives_scope}

    \subsection{Objectives}
    The research questions above translate into four concrete objectives:

    \begin{enumerate}
        \item[\textbf{O1.}] Model each learner's vocabulary word by word, holding a spaced-repetition card for every tracked word and linking it to the dictionary and to the subtitles of every reel (RQ1).

        \item[\textbf{O2.}] Build a three-stage recommendation engine that filters reels for comprehensibility, prioritises those containing words due for review, and ranks the remainder by similarity to the learner's viewing history (RQ2).

        \item[\textbf{O3.}] Generate two vocabulary-driven games, a Wordle-style guessing game and a procedurally generated crossword, from the same vocabulary state, so that every round of play is also a spaced-repetition review (RQ3).
        
        \item[\textbf{O4.}] Deliver these components as a working, offline-capable prototype for Android and iOS, and evaluate their algorithmic behaviour on real seeded content.
    \end{enumerate}

    \subsection{Scope and Delimitations}
    The work is deliberately bounded. The system's own learning mechanics address vocabulary acquisition only. Grammar is not taught or tested within \textit{Glosy}; it is supported only by a curated collection of external YouTube lessons, kept per learning language, which the learner can watch from the application's Videos screen or open in YouTube. The application offers no pronunciation or speaking exercises. It supports three languages --- English, Greek, and Farsi, the last a right-to-left script --- and targets both Android and iOS from one codebase. The two platforms are at different stages of distribution: the Android build is published as versioned APK releases on GitHub, whereas the iOS build is currently available only through Expo. The thesis also treats content supply as outside its scope: the reel catalogue is small and was seeded by a maintainer to demonstrate the recommender, not curated as a learning library (Section~\ref{sec:limitations}). Most importantly, the thesis evaluates the system's algorithmic components, not its effect on learners. No controlled study with learners has been conducted, so no claim is made about vocabulary gain, retention, or engagement, and the content-based ranking stage of the recommender, although implemented, is not tested (Chapter~\ref{chapter:evaluation_results}).

    \subsection{Methodological Approach}
    The thesis follows a design-and-build approach. The literature review (Chapter~\ref{chapter:related_work}) yields the design principles: among them the 98\% lexical-coverage criterion for comprehensible input and a spaced-repetition model of memory decay. Chapters~\ref{methodology} and~\ref{chapter:discussion} turn these principles into an architecture and a working implementation, and Chapter~\ref{chapter:evaluation_results} tests the result. Because no learner has yet used the system, the evaluation examines each component directly. The crossword generator is benchmarked over 4{,}200 trials on word pools sampled from the seeded dictionaries, and the first two recommender stages are run on the real reel catalogue. The benchmark harness and evaluation scripts are part of the project repository, so every reported result can be reproduced (Section~\ref{sec:eval_approach}).

% ----------SIGNIFICANCE----------
\section{Significance}
    The significance of this work lies in three areas.

    For learners, \textit{Glosy} turns time that is already spent on short videos and casual games into vocabulary practice, so that learning does not require finding separate study time --- the workload barrier identified above. It is designed mobile-first, for the device on which those habits already take place, and it supports offline gameplay and practice, which serves learners who study where internet access is limited, a use case that market analyses identify for offline learning features \cite{businessresearch2026languageapps}.

    For research, it operationalises comprehensible input as a measurable criterion: the recommender keeps only reels in which the learner already knows at least 98\% of the words, which turns a largely qualitative construct into a machine-checkable one (Section~\ref{subsec:comprehensible_input}). It also links a single word-level memory model to both content selection and play, so that watching a reel and completing a game round each act as a spaced-repetition review.

    For practice, it is a working prototype rather than a design proposal, and its architecture is language-agnostic: adding a language requires content ingestion, not code changes (Chapter~\ref{chapter:conclusions}). It currently supports English, Greek, and Farsi, including a right-to-left script.

    Because no controlled study with learners has yet been conducted (Chapter~\ref{chapter:evaluation_results}), these are benefits the design is built to deliver rather than effects it has been shown to produce.

\section{Thesis Structure}
The remainder of this thesis is organised as follows. Chapter~\ref{chapter:related_work} reviews mobile-assisted language learning platforms and their limitations, and establishes the pedagogical, memory-modelling, and recommendation foundations from which the design follows. Chapter~\ref{methodology} presents the system architecture, technology stack, data model, and the design of the recommendation engine and the two games. Chapter~\ref{chapter:discussion} documents the implementation of each component --- database schema, backend API, frontend, reels service, and game modules --- with reference to the source code. Chapter~\ref{chapter:evaluation_results} evaluates the system's algorithmic components, namely the Word of Wonders level generator and the recommendation engine, since no controlled study with learners has yet been conducted. Chapter~\ref{chapter:conclusions} states the contributions, acknowledges limitations, and proposes future work.

\newpage